\clearpage
\fancyhead[L]{\small CARLINI. EFFEMERIDI ASTRONOMICHE.}
\hypertarget{native-571-lower}{}
\begin{center}\small Göttingische gelehrte Anzeigen. Stück 48. Seiten 475..478. 24. März 1814.\\[3pt]\rule{.72\linewidth}{.3pt}\end{center}

{\rightskip=0pt plus 3em \emph{Effemeridi astronomiche di Milano per l'anno} 1814, \emph{calcolate da} \textsc{Francesco Carlini}. \emph{Dalla stamperia reale. Con appendice} 120 und 140 Seiten im Kleinquartformat.\par}

Da wir die zweckmäßige Einrichtung dieser vortrefflichen astronomischen Ephemeriden und die ihnen beigegebenen \emph{ständigen} Artikel bereits bei der Besprechung der früheren Jahrgänge beschrieben haben (siehe Stück 95 vom vorigen Jahr [S. 556 d.\,B.]), beschränken wir uns jetzt auf die Besprechung des Anhangs, der auch in diesem Jahrgang sehr reich ausgestattet ist. Zuerst finden wir die Beobachtungen des ersten Kometen des Jahres 1811 von \textsc{Barnabas Oriani}. Sie reichen vom 29. August 1811 bis zum 21. Januar 1812. Mitgeteilt werden sowohl die ursprünglichen Vergleiche am Äquatorsektor als auch die daraus abgeleiteten Positionen, außerdem parabolische Elemente, die diese Beobachtungen sehr gut, die älteren Beobachtungen \textsc{von Zach}s aber weniger genau darstellen. Was \textsc{Oriani} über das äußere Aussehen dieses bemerkenswerten Kometen sagt, stimmt mit den Beobachtungen anderer Astronomen überein. Er vergleicht ihn auch mit dem Kometen von 1744, der in dieser Hinsicht viel Ähnlichkeit mit dem von 1811 hatte. Doch bemerkt dieser einsichtsvolle Astronom mit Recht, dass die Zahl der mit Fernrohren beobachteten größeren Kometen noch viel zu klein, uns es daher zu heikel sei, auf eine gewisse Ähnlichkeit im äußeren Aussehen von drei oder vier Kometen Theorien gründen zu wollen, die den Ursprung der verschiedenen seltsamen Erscheinungen ihrer Atmosphären und Schweife auf plausible Weise erklären könnten. Am 24. Dezember glaubte \textsc{Oriani}, \emph{mitten} im Nebel des Kometenkopfes zum ersten Mal einen hellen Kern zu bemerken. Da er sich in den folgenden Tagen nicht wieder erkennen ließ, wurde Oriani jedoch wieder unsicher und hielt es für wahrscheinlich, dass der vermeintliche Kern ein Fixstern 9. oder 10. Größe gewesen sei. Es scheint uns deshalb wohl der Mühe wert, an der Stelle, an der der Komet am 24. Dezember stand (Rektaszension $306^\circ\,16'\,21''$, Deklination $1^\circ\,54'\,6''$ nördlich), noch einmal nachzusehen (dieser Teil des Himmels ist zur gegenwärtigen Jahreszeit nicht sichtbar). Darauf folgen ebenso ausführlich die Beobachtungen des zweiten Kometen von 1811 durch denselben Astronomen und parabolische Elemente, die allerdings keine sehr gute Übereinstimmung ergeben. Bekanntlich hat Herr \textsc{Nicolai} die Bahn dieses Kometen als Ellipse berechnet. Sie stellt auch die Beobachtungen von \textsc{Oriani} recht gut dar, obwohl sie bei der Berechnung nicht verwendet worden waren. --- Beobachtungen des Kometen von 1812 von \emph{demselben Verfasser} (vom 1. bis zum 25. September) mit parabolischen Elementen. --- Es folgen in einer Reihe von Aufsätzen mehrere beobachtete Planetenpositionen von \textsc{Oriani} und \textsc{Santini}, auf die wir hier nur kurz eingehen können. Von \textsc{Oriani} stammen die Oppositionen des Uranus 1811, des Mars 1811, der Vesta 1811 und der Vesta 1812, die ersten drei am Mauerquadranten von \textsc{Ramsden}, die vierte am Äquatorsektor. Von \textsc{Santini} in Padua stammen die Opposition der Juno 1810 und die des Uranus 1810, beide am Mauerquadranten von \textsc{Ramsden}. --- \SourceLeaf{572}Sternbedeckungen in den Jahren 1812 und 1813 und die Sonnenfinsternis vom 31. Januar 1813, beobachtet von \textsc{Oriani}, mit einer ausführlichen Berechnung der Bedeckung von $\alpha$ im Stier am 22. Oktober 1812 zur Bestimmung der Fehler der neuen Mondtafeln von \textsc{Burckhardt}. Mancher wird sich freuen, die nötigen Vorschriften hier in eleganter Form zusammengestellt und an einem praktischen Beispiel erläutert zu sehen. Vielleicht wäre allerdings für manche Leser der Hinweis nicht ganz überflüssig gewesen, dass gerade in diesem Fall das Ergebnis für den Breitenfehler der Tafeln nicht sehr genau ausfallen konnte. --- Ein Brief von Herrn \textsc{Joseph Piazzi} an \textsc{Oriani} vom 4. Juli 1812 enthält noch einige Verbesserungen seines großen Sternkatalogs und des kleineren im Libro Sesto della specola astronomica. Er bringt außerdem die erfreuliche Nachricht, dass dieser berühmte Beobachter eine ganz neue Bearbeitung des großen Katalogs bereits fast fertiggestellt hat, die die Astronomen nur mit Ungeduld erwarten können. --- Betrachtungen über die astronomischen Uhren von \textsc{Angelo Cesaris} enthalten neben dem Bekannten einige nicht uninteressante Bemerkungen über verschiedene Änderungen an der Hemmung und an anderen Teilen, die italienische Handwerker erprobt haben. Man würde sich nur wünschen, dass diese Bemerkungen etwas ausführlicher wären. --- Noch einige beobachtete und berechnete Oppositionen der neuen Planeten von \textsc{Francesco Carlini}, nämlich Beobachtungen der Ceres 1811 und 1812 und der Pallas 1811 sowie Berechnungen der Oppositionen der Ceres von 1809, 1811 und 1812, der Vesta 1810, der Juno 1810 und der Pallas 1811. Alle diese Ergebnisse wie auch die oben genannten sind für die Planetentheorie von sehr hohem Wert, da sie von sehr geübten Beobachtern stammen und auf Beobachtungen mit vortrefflichen Instrumenten beruhen. --- Den Schluss bilden die dreijährigen meteorologischen Beobachtungen von \textsc{Cesaris} an der Mailänder Sternwarte (1810---1812). Wir erlauben uns dazu hier nur die eine Bemerkung, dass der in den nördlichen Gegenden Europas so ungewöhnlich heiße Sommer von 1811 in Mailand bei Weitem weniger heiß war. In der zweiten Julihälfte betrug der höchste Thermometerstand nur $25^\circ\,4$, und plötzliche Änderungen gab es überhaupt nicht. In Göttingen dagegen stieg das Thermometer im Schatten zum Beispiel am 19. Juli auf $27^\circ$ und am 20. sogar auf $29^\circ$, um dann auf einmal sehr tief zu sinken. Eine niederschlagende Bemerkung für die Liebhaber allgemeiner Wettertheorien!
